% CV of Jui-Te "Ray" Huang
% Build: pixi run cv  (writes assets/pdf/JuiTe_cv.pdf, the PDF the site serves)
\documentclass[11pt]{article}

\usepackage[letterpaper, margin=0.65in]{geometry}
\usepackage[T1]{fontenc}
\usepackage{newtxtext}
\usepackage{enumitem}
\usepackage{titlesec}
\usepackage{xcolor}


\pagestyle{empty}
\setlength{\parindent}{0pt}
\setlength{\parskip}{0pt}

% Section headings: bold, underlined with a full-width rule
\titleformat{\section}{\large\bfseries}{}{0pt}{}[\vspace{-0.6em}\rule{\textwidth}{0.8pt}]
\titlespacing*{\section}{0pt}{1.0em}{0.5em}

% Bullet lists
\setlist[itemize]{leftmargin=1.4em, itemsep=0pt, topsep=1pt, parsep=0pt}
\setlist[itemize,1]{label=\textbullet}
\setlist[itemize,2]{label=$\circ$, leftmargin=1.2em}

% \entry{left bold}{right bold}
\newcommand{\entry}[2]{\textbf{#1}\hfill\textbf{#2}\\}
% \me{name} highlights my name in publication lists
\newcommand{\me}[1]{\textbf{#1}}

\begin{document}

%-------------------------------------------------------------------- Header
\begin{center}
  {\LARGE\bfseries Jui-Te `Ray' Huang}
\end{center}
\vspace{-0.3em}
Webpage: https://huangjuite.github.io
\hfill
Email: juiteh@andrew.cmu.edu

%-------------------------------------------------------------------- Education
\section{Education}
\entry{Carnegie Mellon University}{Pittsburgh, PA}
PhD in Robotics \hfill \textbf{Current}\\
MS in Robotics (MSR) \hfill \textbf{May 2024}

\medskip
\entry{National Yang Ming Chiao Tung University}{Hsinchu, Taiwan}
BS, MS in Electrical and Computer Engineering \hfill \textbf{Aug.\ 2021}

%-------------------------------------------------------------------- Experience
\section{Experience}
\entry{Research Assistant, Carnegie Mellon University}{Sep.\ 2022 -- present, Pittsburgh, PA}
Robot Perception Lab, advised by Prof.\ Michael Kaess.
\begin{itemize}
  \item Radar-Inertial State Estimation and Mapping using mmWave Imaging Radar.
  \begin{itemize}
    \item Research on robust state estimation in adverse environmental conditions.
    \item Research on learning-based radar signal processing for state estimation.
    \item Implemented radar-visual-inertial sensor fusion state estimation.
    \item Developed a large-scale outdoor SLAM system using multiple high-resolution imaging radars.
  \end{itemize}
  \item GPU-accelerated radar signal processing.
  \begin{itemize}
    \item Implemented a MIMO signal processing pipeline that generates radar point clouds from raw digital signals.
    \item Implemented Synthetic Aperture Radar algorithms to create mmWave radar occupancy maps.
  \end{itemize}
  \item Built a multi-modal sensor rig (radar, LiDAR, stereo RGB/thermal/event cameras) and co-led a large-scale data collection campaign.
\end{itemize}

\medskip
\entry{Research Intern, Bosch Research USA}{May 2026 -- Aug.\ 2026, Pittsburgh, PA}
\mbox{}\hfill\textbf{May 2025 -- Aug.\ 2025, Pittsburgh, PA}
\begin{itemize}
  \item Developed a conditional latent diffusion model that synthesizes complex-valued, multi-antenna
    range--Doppler radar spectra from lidar, camera, ego-motion, and radar waveform parameters.
  \item Used a phase-preserving spectrum autoencoder with a flow-matching diffusion transformer, so
    generated spectra can be processed with standard radar signal processing and pretrained networks.
  \item Generated raw radar spectra from CARLA simulator scenes; augmenting radar foundation model
    training with generated data gave gains on the I/Q-1M dataset equivalent to $2\times$ data scaling.
  \item Built a large-scale training and evaluation pipeline; paper in submission.
\end{itemize}

\medskip
\entry{Research Intern, University of Washington}{Sep.\ 2021 -- Feb.\ 2022, Seattle, WA}
Information Processing Lab, advised by Prof.\ Jenq-Neng Hwang.
\begin{itemize}
  \item Built and calibrated a sensor system to collect synchronized data for radar multi-object tracking.
  \item Learning-based multi-object detection and tracking using mmWave imaging radar.
  \item 3D multi-object tracking using LiDAR and visual depth estimation.
\end{itemize}

\medskip
\entry{Research Assistant, National Chiao Tung University}{Jun.\ 2018 -- Aug.\ 2021, Hsinchu, Taiwan}
Assistive Robotics Group, advised by Prof.\ Hsueh-Cheng Wang.
\begin{itemize}
  \item DARPA Subterranean Challenge, SLAM and Deep Reinforcement Learning Navigation.
  \begin{itemize}
    \item Research on navigation using sim-to-real deep reinforcement learning.
    \item Research on cross-modal learning to navigate robots through smoke using mmWave radar.
    \item Team NCTU, Software Team Lead and Human Supervisor in SubT Challenge: Urban Circuit.
    \item Implemented a LiDAR SLAM system for large-scale underground environments.
  \end{itemize}
  \item Maritime RobotX Challenge, Autonomous Surface Vehicle.
  \begin{itemize}
    \item Built the sensor and computation system for an autonomous boat (WAM-V).
  \end{itemize}
\end{itemize}

%-------------------------------------------------------------------- Interests
\section{Research Interests}
Radar Perception, SLAM, State Estimation, Generative Sensor Simulation, Multi-Modal Learning, Field Robotics

%-------------------------------------------------------------------- Publications
\section{Selected Projects \& Publications}
\begin{itemize}[itemsep=4pt]
  \item \textbf{Radar-Inertial Odometry and Mapping}
  \begin{itemize}[itemsep=3pt]
    \item \me{Jui-Te Huang}, Tianshu Huang, Anthony Rowe, and Michael Kaess,
      ``UNRIO: Uncertainty-Aware Velocity Learning for Radar-Inertial Odometry,''
      in submission, 2026. \textbf{Best Student Poster}, ICRA 2026 Radar in Robotics Workshop.
    \item \me{Jui-Te Huang*}, Ruoyang Xu*.
      \"Orebro 4D-Radar SLAM Challenge, \textbf{1st place}. ICRA 2026 Radar in Robotics Workshop.
    \item \me{Jui-Te Huang}, Ruoyang Xu, and Michael Kaess,
      ``Millimeter Wave Radar: From Synthetic Aperture to Probabilistic Mapping,''
      \textit{IEEE/RSJ International Conference on Intelligent Robots and Systems (IROS)}, 2026.
    \item \me{Jui-Te Huang}, Ruoyang Xu, Akshay Hinduja, and Michael Kaess,
      ``Multi-Radar Inertial Odometry for 3D State Estimation using mmWave Imaging Radar,''
      \textit{IEEE International Conference on Robotics and Automation (ICRA)}, 2024.
  \end{itemize}

  \item \textbf{Radar Multi-Object Tracking for Autonomous Driving}
  \begin{itemize}
    \item Y.~Wang, J.-H.~Cheng, \me{J.-T.~Huang}, S.-Y.~Kuan, Q.~Fu, C.~Ni, S.~Hao, G.~Wang,
      G.~Xing, H.~Liu, and J.-N.~Hwang,
      ``Vision meets mmWave Radar: 3D Object Perception Benchmark for Autonomous Driving,''
      \textit{IEEE Intelligent Vehicles Symposium (IV)}, 2023.
  \end{itemize}

  \item \textbf{Search and Rescue in DARPA Subterranean Challenge}
  \begin{itemize}
    \item C.-L.~Lu*, \me{J.-T.~Huang*}, et al.,
      ``A Heterogeneous Unmanned Ground Vehicle and Blimp Robot Team for Search and Rescue
      using Data-driven Autonomy and Communication-aware Navigation,''
      \textit{Field Robotics, Special Issue: Advancements and Lessons Learned during Phase I \& II
      of the DARPA Subterranean Challenge}, 2022. (*Equal contribution)
  \end{itemize}

  \item \textbf{Cross-modal Learning and Reinforcement Learning for Navigation}
  \begin{itemize}[itemsep=3pt]
    \item \me{J.-T.~Huang}, C.-L.~Lu, P.-K.~Chang, C.-I~Huang, C.-C.~Hsu, Z.~L.~Ewe, P.-J.~Huang,
      and H.-C.~Wang,
      ``Cross-Modal Contrastive Learning of Representations for Navigation using Lightweight,
      Low-Cost Millimeter Wave Radar for Adverse Environmental Conditions,''
      \textit{IEEE Robotics and Automation Letters (RA-L)}, 6(2), 3333--3340, 2021.
    \item C.-L.~Lu, Z.-Y.~Liu, \me{J.-T.~Huang}, C.-I~Huang, B.-H.~Wang, Y.~Chen, N.-H.~Wu,
      H.-C.~Wang, L.~Giarr\'e, and P.-Y.~Kuo,
      ``Assistive Navigation using Deep Reinforcement Learning Guiding Robot with UWB/Voice
      Beacons and Semantic Feedbacks for Blind and Visually Impaired People,''
      \textit{Frontiers in Robotics and AI}, 2021.
  \end{itemize}
\end{itemize}

%-------------------------------------------------------------------- Service
\section{Academic Service}
\begin{itemize}
  \item Teaching Assistant, CMU: 16-833 Robot Localization and Mapping, 16-822 Geometry-based Methods in Vision.
  \item Teaching Assistant, NCTU: Sensing and Intelligent Systems, Human-Centric Computing Lab,
    Introduction to Artificial Intelligence.
  \item Reviewer for IEEE IROS, ICRA, RA-L, T-ASE.
\end{itemize}

%-------------------------------------------------------------------- Skills
\section{Skills}
\begin{itemize}
  \item \textbf{Programming Languages:} C++, Python, CUDA
  \item \textbf{Software, Middleware, and Libraries:} ROS, Docker, OpenCV, Open3D, PCL
  \item \textbf{Deep Learning / Optimization Frameworks:} PyTorch, JAX, GTSAM
  \item \textbf{Manufacturing:} Autodesk Fusion, KiCAD, 3D Printing, PCB Design
\end{itemize}

\end{document}
